%=============================================================================
% ALPHA--NEUTRINO SECTOR CROSS-REFERENCE ANALYSIS
% Deep investigation: How the alpha formula connects to neutrino masses
% Generated: 23 March 2026 | Model: Opus 4.6 (1M context)
%=============================================================================

\documentclass[11pt,a4paper]{article}
\usepackage[margin=2.5cm]{geometry}
\usepackage{amsmath,amssymb,amsthm}
\usepackage{booktabs}
\usepackage{enumitem}
\usepackage{float}
\usepackage{xcolor}
\usepackage[most]{tcolorbox}
\usepackage{hyperref}

\newtheorem{theorem}{Theorem}
\newtheorem{lemma}[theorem]{Lemma}

\title{\textbf{Cross-Reference Analysis:\\
How the $\alpha$ Formula Connects to the Neutrino Sector}\\[0.5em]
\large Five Questions on $\alpha$--Neutrino Coupling in DFD}

\author{DFD Research Programme --- Cross-Reference Series (Opus 4.6)}
\date{23 March 2026}

\begin{document}
\maketitle

%=============================================================================
\section*{Executive Summary}
%=============================================================================

\begin{tcolorbox}[colback=yellow!5!white, colframe=red!70!black,
  title=\textbf{KEY FINDINGS}]

\begin{enumerate}
\item \textbf{The charged fermion formula $m_f = A_f \alpha^{n_f} v/\sqrt{2}$
  does NOT apply directly to neutrinos.} Neutrinos use a structurally
  different mechanism: the seesaw with $M_R = M_P \alpha^3$ sets the
  overall scale, and microsector channel fractions (not geodesic-distance
  exponents) set the mass ratios. However, $\alpha$ enters at every stage.

\item \textbf{$M_R = M_P \alpha^3$ is set by the same determinant-scaling
  mechanism as $\alpha^{57}$.} The exponent 3 equals $N_{\rm gen}$, the
  topological generation count from the Atiyah--Singer index on
  $\mathbb{CP}^2 \times S^3$. This is the direct $\alpha$-formula
  connection: the Majorana scale is the Planck mass suppressed by three
  powers of $\alpha$.

\item \textbf{$\Sigma m_\nu = 61.4$~meV vs DESI $< 53$~meV (FC):} The
  tension is model-dependent. Under $w_0 w_a$CDM the bound relaxes to
  $< 163$~meV. DFD's modified cosmology self-consistently relaxes the
  bound. However, no continuous knob in the $\alpha$ formula can reduce
  $\Sigma m_\nu$---the prediction is rigid. CMB-S4 + DESI at
  $\sigma(\Sigma m_\nu) \sim 15$~meV will be decisive.

\item \textbf{$\alpha$ does NOT enter the TBM-to-PMNS corrections
  directly.} The leading correction is $\theta_{13} \neq 0$, which TBM
  sets to zero. DFD does not yet derive $\theta_{13}$ from $\alpha$; it is
  treated as a perturbation using the measured value $s_{13}^2 = 0.022$.
  This is an open problem.

\item \textbf{Normal ordering follows from $r = \alpha^{-7/20} > 1$.}
  Since $\alpha < 1$, any negative exponent gives $r > 1$, meaning
  $m_3 > m_2 > m_1$ (normal ordering). The exponent $7/20 = \dim M / (4n)$
  is set by the topology of $\mathbb{CP}^2 \times S^3$ ($\dim M = 7$) and
  the bundle rank ($n = 5$). Normal ordering is thus a topological
  consequence.
\end{enumerate}
\end{tcolorbox}


%=============================================================================
\part{Question 1: Neutrino Mass Formula vs Charged Fermion Formula}
%=============================================================================

\section{The Charged Fermion Formula}

For the nine charged fermions, DFD predicts (Appendix K, Section 13):
\begin{equation}
\boxed{m_f = A_f \cdot \alpha^{n_f} \cdot \frac{v}{\sqrt{2}}}
\end{equation}
where $n_f \in \{0, 1.0, 1.5, 2.5\}$ are sector-dependent exponents from
$\mathbb{CP}^2$ geodesic distance, and $A_f$ are rational prefactors from
gauge Clebsch--Gordan coefficients, $A_5$ class factors, and QCD running.

\section{Why Neutrinos Are Different}

Neutrinos are fundamentally different because:
\begin{enumerate}
\item They are \textbf{gauge singlets} under $SU(3)_c \times U(1)_{\rm em}$.
  The $\mathbb{CP}^2$ geodesic-distance mechanism that generates the
  exponents $n_f$ for charged fermions operates through gauge coupling
  paths. Neutrinos, being singlets, do not participate in this mechanism
  directly.

\item They acquire mass through the \textbf{Type I seesaw}, not through
  direct Yukawa coupling:
  \begin{equation}
  m_\nu \sim \frac{m_D^2}{M_R}
  \end{equation}
  The Dirac mass $m_D$ connects to the charged-fermion formula, but $M_R$
  introduces a separate, purely microsector-derived scale.

\item The mass \textbf{ratios} are fixed not by geodesic distances but by
  \textbf{discrete symmetry breaking} ($S_3 \to S_2$) and microsector
  channel fractions.
\end{enumerate}

\section{The Neutrino Mass Formula}

The DFD neutrino mass formula (Appendix X, Branch B) is:
\begin{equation}
\boxed{m_3 = \frac{14}{13}\,\pi\,M_P\,\alpha^{14}}
\end{equation}
with ratios:
\begin{align}
\frac{m_2}{m_1} &= \alpha^{-3/11}, \\
\frac{m_3}{m_2} &= \alpha^{-7/20}.
\end{align}

\paragraph{How $\alpha$ enters.}
Count the powers of $\alpha$ in the absolute scale:
\begin{itemize}
\item $M_R = M_P \alpha^3$ (Majorana scale: 3 powers)
\item $v = M_P \alpha^8 \sqrt{2\pi}$ (Higgs VEV: 8 powers)
\item Seesaw: $m_\nu \sim v^2/M_R \propto M_P \alpha^{16}/M_P \alpha^3
  = M_P \alpha^{13}$
\item The finite-$d$ priming factor $14/13$ and the $\pi$ prefactor lift
  this to $m_3 = (14/13)\pi M_P \alpha^{14}$
\end{itemize}

Thus the neutrino mass scale involves $\alpha^{14}$---far more powers of
$\alpha$ than any charged fermion (maximum $\alpha^{2.5}$ for first-generation
particles). This extreme suppression is why neutrinos are so light.

\paragraph{Structural comparison.}
\begin{center}
\begin{tabular}{lcc}
\toprule
& \textbf{Charged fermions} & \textbf{Neutrinos} \\
\midrule
Formula & $m_f = A_f \alpha^{n_f} v/\sqrt{2}$ &
  $m_3 = (14/13)\pi M_P \alpha^{14}$ \\
Exponent origin & $\mathbb{CP}^2$ geodesic distance &
  Seesaw + microsector priming \\
Ratio mechanism & Sector-dependent $n_f$ &
  $S_3 \to S_2$ channel fractions \\
Free parameters & Zero (given microsector) & Zero (given microsector) \\
$\alpha$ powers & $0$--$2.5$ & $14$ (absolute), fractional (ratios) \\
\bottomrule
\end{tabular}
\end{center}


%=============================================================================
\part{Question 2: What Sets $M_R$? Connection to $\alpha$}
%=============================================================================

\section{The Derivation of $M_R = M_P \alpha^3$}

This is derived in Appendix P (Theorem P.3) using the same mechanism as the
cosmological constant derivation ($\rho_{\rm vac}/\rho_{\rm Pl} = \alpha^{57}$,
Appendix O).

\paragraph{The mechanism.}
Let $\mathcal{H}_{\nu_R}$ be the internal Hilbert space of right-handed
neutrinos. By the Atiyah--Singer index theorem on $\mathbb{CP}^2 \times S^3$:
\begin{equation}
\dim(\mathcal{H}_{\nu_R}) = N_{\rm gen} = 3.
\end{equation}

The singlet-sector Toeplitz operator $\mathcal{K}_{\nu_R}$ has the unique
dimensionless suppression factor:
\begin{equation}
\varepsilon_{\nu_R}(g) = \frac{\det(\mathcal{K}_{\nu_R})}
  {\det(g\,\mathcal{K}_{\nu_R})} = g^{-N_{\rm gen}} = \alpha^{N_{\rm gen}}
  = \alpha^3,
\end{equation}
where $g = \alpha^{-1}$ is the microsector coupling. Hence:
\begin{equation}
\boxed{M_R = M_P\,\alpha^3 \approx 4.74 \times 10^{12}~\text{GeV}}
\end{equation}

\section{Parallel with $\alpha^{57}$}

\begin{center}
\small
\begin{tabular}{lll}
\toprule
& \textbf{Cosmological constant} & \textbf{Majorana scale} \\
\midrule
State space dim & $k_{\max} = 60$ & $N_{\rm gen} = 3$ \\
Kernel dim & $N_{\rm gen} = 3$ & 0 (no kernel) \\
Exponent & $60 - 3 = 57$ & $3$ \\
Result & $\rho_{\rm vac}/\rho_{\rm Pl} = \alpha^{57}$ &
  $M_R/M_P = \alpha^3$ \\
\bottomrule
\end{tabular}
\end{center}

Both exponents are \textbf{topological integers}: 57 from the Spin$^c$ index
minus the generation kernel, 3 from the generation count alone. Neither is
fitted.

\section{$M_R$ and $\alpha$: The Direct Link}

The $\alpha$ formula enters $M_R$ through exactly one mechanism: the
determinant scaling of the Toeplitz-quantized singlet sector with coupling
$g = \alpha^{-1}$. The same $\alpha$ that is derived from the Chern--Simons
level sum at $k_{\max} = 60$ (Appendix K) sets:
\begin{itemize}
\item The electromagnetic coupling strength
\item The Majorana scale (via $\alpha^3$)
\item The electroweak VEV (via $\alpha^8$)
\item The neutrino mass scale (via $\alpha^{14} = \alpha^{2 \times 8 - 3 + 1}$)
\end{itemize}

This is the deepest connection: $\alpha$ is the single universal
dimensionless coupling, and the neutrino sector is controlled by its 14th
power.


%=============================================================================
\part{Question 3: $\Sigma m_\nu = 61.4$~meV vs DESI Bound}
%=============================================================================

\section{The Tension}

\begin{center}
\begin{tabular}{lcc}
\toprule
Constraint & 95\% upper bound & DFD status \\
\midrule
$\Lambda$CDM (DESI DR2+Planck) & $\Sigma m_\nu < 64$~meV & MARGINAL \\
$\Lambda$CDM (DESI+Planck+CMB lensing, FC) & $\Sigma m_\nu < 53$~meV &
  TENSION \\
$w_0 w_a$CDM (DESI+Planck) & $\Sigma m_\nu < 163$~meV & SAFE \\
\bottomrule
\end{tabular}
\end{center}

The Feldman--Cousins corrected bound of 53~meV is below both the DFD
prediction (61.4~meV) and the oscillation-experiment minimum for normal
ordering ($\sim 59$~meV). This means the 53~meV bound, if taken at face
value, would exclude \emph{all} normal-ordering models with standard
cosmology.

\section{Does the $\alpha$ Formula Give Flexibility?}

\textbf{No.} The DFD prediction is rigid:
\begin{enumerate}
\item The mass ratios $k = \alpha^{-3/11}$ and $r = \alpha^{-7/20}$ are
  fixed by locked microsector integers. There is no continuous parameter
  to adjust.

\item The absolute scale $m_3 = (14/13)\pi M_P \alpha^{14}$ depends only
  on:
  \begin{itemize}
  \item $M_P$ (measured)
  \item $\alpha$ (derived, locked at $1/137.036$)
  \item The integers 14 and 13 (from $a + 5$ and $a + 4$ with $a = 9$)
  \end{itemize}
  None of these can be varied.

\item The only discrete freedom is the Branch choice (A vs B). Branch A
  gives $\Sigma m_\nu = 69.2$~meV (worse), so there is no escape in the
  other direction either.
\end{enumerate}

\paragraph{What could change.}
If the minimal-padding prescription for $(a, n)$ were modified---for
instance if a different bundle choice gave $(a, n) = (10, 5)$---then
$d_\nu = a + 5 = 15$, $F_\nu = 15/14$, and the masses would shift. But
$(a, n) = (9, 5)$ is forced by the same construction that gives
$k_{\max} = 60$ and $\alpha = 1/137$. Changing $(a, n)$ would destroy the
$\alpha$ prediction.

\section{Self-Consistency Argument}

DFD's cosmology differs from $\Lambda$CDM: the psi-screen modifies the
distance--redshift relation, and $G_{\rm eff}(k)$ modifies growth. The
neutrino mass bound is extracted \emph{assuming} $\Lambda$CDM. Under DFD
cosmology:
\begin{itemize}
\item The CMB lensing power spectrum changes (modified growth)
\item The BAO scale interpretation changes (psi-screen optical bias)
\item The resulting bound relaxes to $\sim 100$--$163$~meV
\end{itemize}

This is a valid self-consistency argument but has the weakness that it is
unfalsifiable until a DFD Boltzmann code exists. The urgent need is:
\begin{equation}
\text{DFD-modified CLASS/CAMB} \longrightarrow
\text{DFD-consistent } \Sigma m_\nu \text{ bound}
\end{equation}

\section{Timeline for Resolution}

\begin{center}
\begin{tabular}{lll}
\toprule
Experiment & Sensitivity & Date \\
\midrule
JUNO & Mass ordering (pass/fail) & 2027 \\
DESI + CMB-S4 & $\sigma(\Sigma m_\nu) \sim 15$~meV & 2027--2029 \\
Euclid + ground lensing & $\sigma(\Sigma m_\nu) \sim 20$~meV & 2028 \\
nEXO & $m_{\beta\beta}$ to 5~meV & 2030+ \\
\bottomrule
\end{tabular}
\end{center}

At $\sigma = 15$~meV, a measurement of $\Sigma m_\nu = 61 \pm 15$~meV
would be a $4\sigma$ detection and a triumph for DFD. A measurement of
$\Sigma m_\nu < 30$~meV (even model-independently) would falsify it.


%=============================================================================
\part{Question 4: Does $\alpha$ Enter the PMNS Corrections?}
%=============================================================================

\section{TBM as the Leading-Order PMNS}

DFD derives tribimaximal (TBM) mixing from the ``neutrinos-at-center''
overlap rule on $\mathbb{CP}^2$:
\begin{equation}
U_{\rm TBM} = \begin{pmatrix}
\sqrt{2/3} & \sqrt{1/3} & 0 \\
-\sqrt{1/6} & \sqrt{1/3} & \sqrt{1/2} \\
\sqrt{1/6} & -\sqrt{1/3} & \sqrt{1/2}
\end{pmatrix}
\end{equation}

This gives $\theta_{12} = \arcsin(1/\sqrt{3}) \approx 35.3^\circ$,
$\theta_{23} = 45^\circ$, $\theta_{13} = 0$, and $\delta_{CP} = 0$.

\section{The $\theta_{13} \neq 0$ Problem}

Experiment measures $s_{13}^2 = 0.0220 \pm 0.0007$. TBM predicts
$s_{13}^2 = 0$. This is the single largest discrepancy in the DFD
neutrino sector.

\paragraph{Does $\alpha$ enter the correction?}
In principle, if $\theta_{13}$ arose from a higher-order perturbation on
$\mathbb{CP}^2$, it should scale as some power of $\alpha$. Note:
\begin{equation}
s_{13}^2 = 0.022 \approx 3\alpha \approx \frac{3}{137} = 0.0219.
\end{equation}

This numerical coincidence is suggestive but \textbf{not derived} in the
current DFD framework. The appendices treat $\theta_{13}$ as an external
perturbation (using the measured value to compute $m_\beta$ and
$m_{\beta\beta}$ with realistic phases), not as a predicted quantity.

\paragraph{What would be needed.}
A derivation of $\theta_{13}$ from the microsector would require:
\begin{enumerate}
\item A mechanism to perturb TBM (e.g., RG running from $M_R$ to
  low scales, or subleading $\mathbb{CP}^2$ overlap corrections)
\item The perturbation to give $s_{13}^2 \propto \alpha^p$ for some $p$
\item The Dirac CP phase $\delta_{CP}$ to emerge simultaneously
\end{enumerate}

This remains an \textbf{open problem} in DFD. The $s_{13}^2 \approx 3\alpha$
observation is noted as a potential clue but is not yet a derivation.

\section{Status of PMNS Corrections}

\begin{center}
\begin{tabular}{lccc}
\toprule
Parameter & TBM (DFD) & Measured & $\alpha$-derived? \\
\midrule
$\sin^2\theta_{12}$ & $1/3 = 0.333$ & $0.307 \pm 0.013$ & No (2$\sigma$) \\
$\sin^2\theta_{23}$ & $1/2 = 0.500$ & $0.546 \pm 0.021$ & No (2$\sigma$) \\
$\sin^2\theta_{13}$ & $0$ & $0.0220 \pm 0.0007$ & No (open) \\
$\delta_{CP}$ & $0$ & $\sim 195^\circ$ & No (open) \\
\bottomrule
\end{tabular}
\end{center}

\textbf{Bottom line:} $\alpha$ does not currently enter the PMNS
corrections in DFD. The corrections from TBM to the physical PMNS matrix
are an open problem. The framework's strength is in the mass eigenvalues
(where $\alpha$ enters through the ratios and absolute scale), not in the
mixing angles beyond leading TBM.


%=============================================================================
\part{Question 5: Normal Ordering from Topology}
%=============================================================================

\section{How Normal Ordering Arises}

The neutrino mass hierarchy parameter is:
\begin{equation}
r = \frac{m_3}{m_2} = \alpha^{-7/20}
\end{equation}

Since $\alpha \approx 1/137 < 1$ and the exponent is $-7/20 < 0$:
\begin{equation}
r = \alpha^{-7/20} = (137.036)^{7/20} \approx 5.60 > 1
\end{equation}

Therefore $m_3 > m_2$. Combined with $k = m_2/m_1 > 1$ (also from a
negative $\alpha$-exponent), we have:
\begin{equation}
m_3 > m_2 > m_1 \qquad \Longleftrightarrow \qquad \text{Normal Ordering}
\end{equation}

\section{The Topological Origin of $r$}

The exponent $7/20$ comes from:
\begin{equation}
\frac{7}{20} = \frac{\dim M}{4n}
\end{equation}
where:
\begin{itemize}
\item $\dim M = \dim(\mathbb{CP}^2 \times S^3) = 4 + 3 = 7$ (real
  dimension of the internal manifold)
\item $n = 5$ is the trivial-bundle rank in $E = \mathcal{O}(9) \oplus
  \mathcal{O}^{\oplus 5}$
\item The factor of 4 in the denominator comes from the real dimension
  of $\mathbb{CP}^2$
\end{itemize}

All three integers (7, 4, 5) are topological:
\begin{itemize}
\item 7 is the total internal dimension
\item 4 is $\dim_{\mathbb{R}}(\mathbb{CP}^2)$
\item 5 is the rank of the trivial summand in the twist bundle, fixed by
  the minimal-padding construction that gives $k_{\max} = 60$
\end{itemize}

\paragraph{Could inverted ordering arise?}
Only if $r < 1$, which requires the exponent $\dim M/(4n)$ to be
\emph{positive} in the $\alpha^{+\cdot}$ convention. Since $\alpha < 1$
and the hierarchy is a \emph{suppression} (the singlet is heavier than
the doublet, a universal feature of seesaw models), the sign is fixed.
Inverted ordering would require either:
\begin{enumerate}
\item A different internal manifold with different dimension ratios, or
\item A fundamentally different mass-generation mechanism for neutrinos
\end{enumerate}

Neither is available within DFD. Normal ordering is a \textbf{topological
prediction}.

\section{JUNO as the Decisive Test}

JUNO (Jiangmen Underground Neutrino Observatory) will determine the mass
ordering at $> 3\sigma$ by $\sim$2027 through reactor antineutrino
oscillations. The DFD prediction is:
\begin{equation}
\boxed{\text{Normal Ordering (NO)}}
\end{equation}

If JUNO confirms inverted ordering, DFD's neutrino sector is falsified.
This is explicitly listed as a falsification criterion in the PRL
supplementary (``Neutrino mass ordering is inverted'' $\Rightarrow$
falsified).


%=============================================================================
\part{Synthesis: The $\alpha$--Neutrino Connection Map}
%=============================================================================

\section{Complete Dependency Graph}

The $\alpha$ formula connects to the neutrino sector through the following
chain:

\begin{tcolorbox}[colback=green!5!white, colframe=green!60!black,
  title=\textbf{$\alpha$ $\to$ Neutrino Connection Map}]

\textbf{Level 0: Topology}
\begin{itemize}[nosep]
\item $\mathbb{CP}^2 \times S^3$, twist bundle
  $E = \mathcal{O}(9) \oplus \mathcal{O}^{\oplus 5}$
\item $k_{\max} = \chi(\mathbb{CP}^2, E) = 60$
\item $N_{\rm gen} = 3$ (Atiyah--Singer index)
\item $\dim M = 7$, $(a, n) = (9, 5)$
\end{itemize}

\textbf{Level 1: $\alpha$ derivation}
\begin{itemize}[nosep]
\item Chern--Simons level sum at $k_{\max} = 60$ $\Rightarrow$
  $\alpha^{-1} = 137.036$
\end{itemize}

\textbf{Level 2: Electroweak scale}
\begin{itemize}[nosep]
\item $v = M_P \alpha^8 \sqrt{2\pi} = 246$~GeV ($\alpha$ enters with
  power 8)
\end{itemize}

\textbf{Level 3: Majorana scale}
\begin{itemize}[nosep]
\item $M_R = M_P \alpha^3 = 4.74 \times 10^{12}$~GeV ($\alpha$ enters
  with power 3)
\end{itemize}

\textbf{Level 4: Seesaw scale}
\begin{itemize}[nosep]
\item $m_\nu \sim v^2/M_R \sim M_P \alpha^{13}$ ($\alpha$ enters with
  power 13)
\end{itemize}

\textbf{Level 5: Mass ratios}
\begin{itemize}[nosep]
\item $k = \alpha^{-3/11}$: from $K^{-1}$ degree / $(\dim M + 4)$
\item $r = \alpha^{-7/20}$: from $\dim M / (4n)$
\item Both from locked microsector integers
\end{itemize}

\textbf{Level 6: Absolute scale closure}
\begin{itemize}[nosep]
\item $d_\nu = a + 5 = 14$, $F_\nu = 14/13$
\item $m_3 = (14/13)\pi M_P \alpha^{14}$
\end{itemize}

\textbf{Level 7: Predictions}
\begin{itemize}[nosep]
\item $(m_1, m_2, m_3) = (2.34, 8.96, 50.12)$~meV
\item $\Sigma m_\nu = 61.4$~meV
\item $\chi^2 = 0.025$ vs NuFIT 6.0 (2 dof, $p = 0.99$)
\end{itemize}
\end{tcolorbox}

\section{Striking Arithmetic Identity}

The combined hierarchy exponent for Branch B contains $\alpha^{-1} = 137$
as an arithmetic consequence:
\begin{equation}
k^2 r^2 = \alpha^{-(6/11 + 7/10)} = \alpha^{-137/110}
\end{equation}

The numerator 137 is the inverse fine-structure constant. This is not
fitted---it emerges from $6/11 + 7/10 = (60 + 77)/110 = 137/110$ using
the locked rational exponents.

\section{Open Problems}

\begin{enumerate}
\item \textbf{First-principles Branch B selection.} Why is Branch B
  ($k = \alpha^{-3/11}$) preferred over Branch A ($k = \alpha^{-2/13}$)?
  Currently selected by data match ($\chi^2 = 0.025$ vs 14.5).

\item \textbf{PMNS corrections from $\alpha$.} Can $\theta_{13}$,
  $\delta_{CP}$ be derived from the microsector? The coincidence
  $s_{13}^2 \approx 3\alpha$ is suggestive.

\item \textbf{DFD Boltzmann code.} Needed to compute the DFD-consistent
  $\Sigma m_\nu$ bound and resolve the DESI tension.

\item \textbf{Absolute scale cross-check.} The exponent 14 in
  $m_3 \propto \alpha^{14}$ equals $2 \times 8 - 3 + 1 = 14$ (two
  Higgs VEV powers minus Majorana power plus priming). Can this be
  derived more transparently?

\item \textbf{Majorana phases.} DFD does not predict the two Majorana
  phases in the PMNS matrix. These affect $m_{\beta\beta}$ predictions
  for neutrinoless double beta decay.
\end{enumerate}


%=============================================================================
\section*{Summary Table: Where $\alpha$ Enters the Neutrino Sector}
%=============================================================================

\begin{table}[H]
\centering
\small
\caption{Every appearance of $\alpha$ in the DFD neutrino sector.}
\begin{tabular}{llc}
\toprule
\textbf{Quantity} & \textbf{Formula} & \textbf{$\alpha$ power} \\
\midrule
Electroweak VEV & $v = M_P \alpha^8 \sqrt{2\pi}$ & 8 \\
Majorana scale & $M_R = M_P \alpha^3$ & 3 \\
Seesaw scale & $m_\nu \sim v^2/M_R$ & $2(8) - 3 = 13$ \\
Primed absolute scale & $m_3 = (14/13)\pi M_P \alpha^{14}$ & 14 \\
Doublet splitting & $m_2/m_1 = \alpha^{-3/11}$ & $-3/11$ \\
Singlet-doublet ratio & $m_3/m_2 = \alpha^{-7/20}$ & $-7/20$ \\
Combined hierarchy & $k^2 r^2 = \alpha^{-137/110}$ & $-137/110$ \\
TBM mixing & Geometry, no $\alpha$ & --- \\
$\theta_{13}$ correction & Not derived (open) & ? \\
\bottomrule
\end{tabular}
\end{table}


\end{document}
